%%%%%%%%%%%%%%%%%%%%%%preface.tex%%%%%%%%%%%%%%%%%%%%%%%%%%%%%%%%%%%%%%%%%
% sample preface
%
% Use this file as a template for your own input.
%
%%%%%%%%%%%%%%%%%%%%%%%% Springer %%%%%%%%%%%%%%%%%%%%%%%%%%

%\preface

\chapter*{Prefacio}
\markboth{Prólogo}{Prefacio}
\addcontentsline{toc}{chapter}{Prefacio}

El presente documento expone los fundamentos teóricos, matemáticos y algorítmicos que sustentan la investigación titulada ``Computabilidad en autómatas celulares complejos''. El propósito central de este trabajo es formalizar y demostrar la capacidad de ciertos sistemas dinámicos discretos, específicamente bajo las reglas de transición \textit{Spiral} y \textit{Beehive}, para ejecutar procesos de cómputo universal a través de la colisión determinista de configuraciones móviles o \textit{gliders}. 

A lo largo del desarrollo de la ingeniería y las ciencias de la computación, es recurrente observar cómo problemas que inicialmente parecen abstractos o puramente teóricos exigen soluciones arquitectónicas y algorítmicas que permiten comprender el manejo de la información a bajo nivel. Bajo esta premisa, el trabajo se encuentra estructurado en una progresión axiomática que guía al lector desde las definiciones más fundamentales hasta el análisis de complejidad algorítmica.

En el Capítulo \ref{cap:preliminares}, se delimita formalmente el objeto de estudio, estableciendo los criterios analíticos y sintéticos necesarios para evaluar la emergencia computacional en los modelos seleccionados. Posteriormente, el Capítulo \ref{cap:teoria_formal} consolida el marco conjuntista y topológico, abordando el concepto crítico de la irreversibilidad en la evolución del autómata; se demuestra cómo la aniquilación de información y la generación de estados Jardín del Edén son requisitos axiomáticos para emular compuertas lógicas tradicionales.

El diseño de un motor matemático computacional exige una comprensión profunda de las estructuras espaciales subyacentes. Por ello, el Capítulo \ref{cap:geometria_computacional} detalla la geometría analítica de la teselación hexagonal y justifica las decisiones de implementación, comprobando teóricamente que la complejidad temporal de evaluación y la actualización mediante \textit{double buffering} mantienen una cota lineal $\mathcal{O}(n)$. Esta base permite formalizar, en el Capítulo \ref{chap:isomorfismos}, el mapeo topológico esencial del proyecto: la transformación biyectiva de una matriz hexagonal a una subred ortogonal de memoria. Mediante el modelo algebraico de transformación definido en la Ecuación \eqref{eq:mapeo_direccional} y la preservación de la métrica de distancia de la Ecuación \eqref{eq:distancia_hexagonal}, se garantiza la invariancia espacial, ilustrada estructuralmente en la Figura \ref{fig:tensor_vecindad_2x2}.

Finalmente, el Capítulo \ref{cap:dinamica_particulas} aborda la fenomenología de estas partículas en movimiento, culminando con la evaluación algorítmica de sus interacciones. Se formaliza que la sincronización determinista de \textit{gliders} para el diseño de circuitos lógicos recae en la clase de complejidad NP, validando la necesidad de los métodos de búsqueda iterativa descritos por la Ecuación \eqref{eq:complejidad_busqueda}.

Este documento está dirigido a estudiantes de ingeniería, matemáticos e investigadores interesados en la computación no convencional, la teoría de autómatas y el análisis de sistemas complejos. La comprensión competente de estos modelos lógicos y espaciales ofrece una perspectiva rigurosa sobre cómo reglas simples, ejecutadas masivamente en paralelo, pueden producir resultados coherentes, estables y computacionalmente universales.

\vspace{\baselineskip}
\begin{flushright}\noindent
Ciudad de México, México\hfill {\it Esparza Jacobo Diego}\\
Estado de México, México\hfill {\it Pérez Velázquez Leonardo Daniel}\\
Abril de 2026\hfill ~\\
\end{flushright}